% Release supplement, 2026-09-16. Shared by all three manual editions.
\clearpage
\fancyhead[R]{\footnotesize\color{GOMuted} MCP Release · 2026-09-16}
\setmonofont{Menlo}[Scale=0.82,Mapping={}]
\newcommand{\McpCode}[1]{\texttt{\detokenize{#1}}}

\CNPage{44}{MCP：连接桌面应用}{MCP: connect the desktop app}{开发版补充 · 2026-09-16。MCP 让外部 AI 客户端调用 ZenOffice 的应用服务。本节描述当前源码功能，不表示已发布安装版具备全部接口。}{
\textbf{准备与启动}
\begin{enumerate}
\item 安装 Node.js 22.12 或更新版本，在仓库根目录运行 \McpCode{npm ci} 和 \McpCode{npm run build:all}。
\item 使用 \McpCode{openssl rand -hex 32} 生成本地令牌。完全退出已运行的 ZenOffice；在同一终端设置下列变量并启动：
\end{enumerate}
\begin{flushleft}\small
\McpCode{export ZENOFFICE_MCP_PORT=19385}\\
\McpCode{export ZENOFFICE_MCP_TOKEN="YOUR_LOCAL_TOKEN"}\\
\McpCode{./node_modules/.bin/electron apps/shell}
\end{flushleft}
将下面配置加入客户端；两处路径改成本机仓库路径，令牌替换为上面生成的同一个值。
\VerbatimInput[fontsize=\small]{mcp-client.example.json}
\Tip{客户端应重新连接并列出工具。连接桌面时应有 38 个工具；不配置 env 时只有 6 个独立可视化工具。调用 \McpCode{application_status} 检查连接，调用 \McpCode{application_list_tabs} 获取实时标签列表。服务默认关闭，仅监听本机；令牌至少 32 字符，不要提交到仓库。}
}

\ENPage{44}{MCP：连接桌面应用}{MCP: connect the desktop app}{Release supplement · 2026-09-16. MCP exposes application services to an external AI client. These instructions describe current source functionality, not a claim about the published installer.}{
\textbf{Prepare and launch}
\begin{enumerate}
\item Install Node.js 22.12 or newer. Run \McpCode{npm ci} and \McpCode{npm run build:all} in the repository root.
\item Generate a local token with \McpCode{openssl rand -hex 32}. Quit any running ZenOffice instance, then launch from the same terminal:
\end{enumerate}
\begin{flushleft}\small
\McpCode{export ZENOFFICE_MCP_PORT=19385}\\
\McpCode{export ZENOFFICE_MCP_TOKEN="YOUR_LOCAL_TOKEN"}\\
\McpCode{./node_modules/.bin/electron apps/shell}
\end{flushleft}
Add this configuration to the client. Replace both paths with your repository path and use the same generated token:
\VerbatimInput[fontsize=\small]{mcp-client.example.json}
\Tip{Reconnect and list tools: expect 38 with the desktop bridge, or 6 independent visualization tools without env. Check \McpCode{application_status}, then fetch live tabs with \McpCode{application_list_tabs}. The bridge is disabled by default and listens only on loopback. Tokens must be at least 32 characters; never commit them.}
}

\CNPage{45}{MCP：选图与插入 PDF}{MCP: choose a chart and insert into PDF}{与 Excel 引导使用相同的图表绑定、AI 建议校验和渲染逻辑。字段必须由实际数据提供，生成结果先预览再插入。}{
\begin{enumerate}
\item 调用 \McpCode{visualization_suggest}：传入 table（columns 列名数组及 rows 行数组）、selectedColumns（允许使用的列名）和 instruction（分析目的）。例如列为“地区、收入、备注”，仅选择“地区、收入”，目的为“比较各地区收入”。
\item 返回 suggestion（图表类型、横轴、显示列、标题、理由）、request（可直接用于渲染）、svg 和 saved=false。只有所选列概况及前 30 行发给模型；完整所选数据用于校验。整个输入先经过本机客户端，敏感列可在调用前移除。
\item 审阅并按需修改 request 的 x、y、title，然后调用 \McpCode{visualization_render_svg}。客户端将返回的 svg 文本保存为本机 \McpCode{/absolute/path/chart.svg}。渲染工具本身不写磁盘。
\item 在应用打开 PDF，再调用 \McpCode{images_targets}，选择 kind=pdf 的最新目标 ID；调用 \McpCode{images_insert}，传入 targetId 和 SVG 的绝对路径 path。
\item 收到 inserted=true 后，PDF 当前页已插入图片并进入编辑状态。可移动、缩放，最后在 PDF 界面保存。saved=false 表示接口未主动保存。
\end{enumerate}
\textbf{本地文件与其他目标}
\begin{itemize}
\item PDF 的“编辑 $\rightarrow$ 插入图片”允许直接选择 .svg，然后点击页面放置。SVG 会转为图片，不保留可编辑矢量路径。
\item PNG/JPEG 可插入 Word、PPT、MD、PDF；MCP 的 SVG 当前仅支持 PDF。单文件上限 15 MB。
\item 新建未保存 MD 可接收内嵌图片；已有路径的 MD 使用旁边的 assets 目录。插入不等于保存。
\item 新开或关闭文件后，重新调用 \McpCode{images_targets}；Excel 分享面板会自动刷新。超时先检查目标，避免重复插入。
\end{itemize}
\Tip{AI 选图需已配置桌面模型，可能产生调用费用；建议客户端超时设为 1200000 毫秒。没有模型时，\McpCode{visualization_profile} 提供本地推荐，\McpCode{visualization_catalog} 提供字段规则。当前支持 81/104 种图形，以目录的可用标记为准。}
}

\ENPage{45}{MCP：选图与插入 PDF}{MCP: choose a chart and insert into PDF}{The workflow shares chart bindings, AI validation, and rendering with the Excel guide. Use actual data fields and review the preview before insertion.}{
\begin{enumerate}
\item Call \McpCode{visualization_suggest} with table (columns and rows), selectedColumns (allowed field names), and instruction (analysis goal). For Region, Revenue, Notes, select only Region and Revenue to compare revenue by region.
\item The result contains suggestion (chart, axes, fields, title, reason), request (render-ready input), svg, and saved=false. Only selected-column profiles and the first 30 rows reach the model; validation uses all selected data. Remove sensitive fields before calling if the local client must not receive them either.
\item Review or change x, y, and title in request, then call \McpCode{visualization_render_svg}. A file-capable client saves the returned svg text as \McpCode{/absolute/path/chart.svg}; the renderer does not write files.
\item Open a PDF, call \McpCode{images_targets}, and choose the latest kind=pdf target ID. Call \McpCode{images_insert} with targetId and the absolute SVG path in path.
\item inserted=true confirms insertion on the current PDF page. Move or resize the image, then save in the PDF UI. saved=false means the tool did not explicitly save.
\end{enumerate}
\textbf{Local files and other targets}
\begin{itemize}
\item PDF Edit $\rightarrow$ Insert Image also accepts .svg files; click the page to place the image. SVG is rasterized; editable vector paths are not retained.
\item PNG/JPEG work in Word, PPT, Markdown, and PDF. SVG through MCP currently requires a PDF target. The file limit is 15 MB.
\item Untitled Markdown accepts embedded images; saved Markdown uses the adjacent assets directory. Insertion does not imply saving.
\item Refresh \McpCode{images_targets} after opening or closing tabs; Excel refreshes its sharing list automatically. Inspect the target before retrying a timeout.
\end{itemize}
\Tip{AI recommendations use the configured desktop model and may incur provider charges. Use a 1200000 ms client timeout. Without AI, use \McpCode{visualization_profile} for local recommendations and \McpCode{visualization_catalog} for bindings. Currently 81 of 104 chart types render; consult availability flags.}
}

\CNPage{46}{MCP：服务清单与写回}{MCP: services and document writes}{MCP 指应用服务接口，不仅是图表服务。当前覆盖如下；未列出的编辑操作不能当作已有工具。}{
\begin{itemize}
\item 可视化 6 项：目录、数据概况、SVG 渲染、仪表盘 SVG、仪表盘 JSON 导出和导入。仪表盘支持 1–6 张图表、1–3 列及共享分类字段筛选。
\item 应用 5 项：状态、标签列表、激活、打开文件、新建文档。
\item 项目与对话 8 项、知识库 6 项：项目列表/创建/重命名、文件与对话读取、知识检索和设置等。
\item 图片共享 2 项：\McpCode{images_targets}、\McpCode{images_insert}。
\item AI 5 项：\McpCode{ai_status}、\McpCode{ai_chat}、\McpCode{visualization_suggest}、\McpCode{ai_review}、\McpCode{ai_screenwriting}。
\item Markdown 3 项：\McpCode{markdown_read}、\McpCode{markdown_replace}、\McpCode{markdown_save}。
\item Word 3 项：\McpCode{word_read_text}、\McpCode{word_insert_text}、\McpCode{word_save}。
\end{itemize}
\textbf{安全写回顺序}
\begin{enumerate}
\item 从标签列表取得 id。Markdown 先 read，再把读取的 text 原样作为 replace 的 expectedText，另传新 text；不一致时重新读取合并。最后 save。
\item Word 先 read\_text，保留 revision；insert\_text 传 id、text、expectedRevision 和 position（start/end）。插入为纯文本段落，可在保存前撤销，再调用 save。
\item 两种保存工具都要求已有文件路径；未命名文档先在界面保存。Word 检测到文件被外部修改时拒绝覆盖。
\end{enumerate}
\Tip{AI 编剧专用入口位于 Word 和 MD，Excel 不提供。审稿与编剧服务见下一节。公文排版、ArtFlow 生图、PDF 保存及全部 Excel/PPT 编辑操作仍未完整接入 MCP。完整参数、示例和错误规则见仓库 \McpCode{docs/mcp-usage.md}；以客户端实时工具列表为准。}
}

\ENPage{46}{MCP：服务清单与写回}{MCP: services and document writes}{MCP covers application services, beyond charts. Current coverage is listed below; unlisted editing operations must not be assumed available.}{
\begin{itemize}
\item Visualization (6): catalog, profile, SVG rendering, dashboard SVG, dashboard JSON export and import. Dashboards support 1–6 cards, 1–3 columns, and a shared categorical filter.
\item Application (5): status, tabs, activation, open file, create document.
\item Projects/chats (8) and knowledge (6): project creation and renaming, file/chat reads, knowledge search and settings.
\item Images (2): \McpCode{images_targets}, \McpCode{images_insert}.
\item AI (5): \McpCode{ai_status}, \McpCode{ai_chat}, \McpCode{visualization_suggest}, \McpCode{ai_review}, \McpCode{ai_screenwriting}.
\item Markdown (3): \McpCode{markdown_read}, \McpCode{markdown_replace}, \McpCode{markdown_save}.
\item Word (3): \McpCode{word_read_text}, \McpCode{word_insert_text}, \McpCode{word_save}.
\end{itemize}
\textbf{Document write sequence}
\begin{enumerate}
\item Get id from the live tab list. Read Markdown first; pass the exact returned text as expectedText to replace, together with the new text. On conflict, read again and merge. Then save.
\item Read Word text and keep revision. Insert text with id, text, expectedRevision, and position (start/end). This adds literal paragraphs and supports undo before saving. Then call save.
\item Both save tools require an existing file path. Save untitled documents in the UI first. Word refuses to overwrite externally changed files.
\end{enumerate}
\Tip{Dedicated screenwriting UI belongs to Word and Markdown, not Excel. Review and screenwriting tools are described next. Official-document formatting, ArtFlow generation, PDF saving, and full Excel/PPT editing remain incomplete in MCP. See \McpCode{docs/mcp-usage.md} for full parameters and errors; use the client's live tool list as the source of truth.}
}

\CNPage{47}{AI 审稿与编剧：统一流程}{AI review and screenwriting: shared workflow}{Word、Markdown 界面与 MCP 复用相同的模型设置、编剧技能和审稿委员会执行流程。先确认素材与任务，再审阅结果，最后写回文档。}{
\textbf{界面操作}
\begin{itemize}
\item AI 审稿：选择标准、报告语言及是否检索文献；查看委员独立意见、主席汇总和失败信息。报告不会自动替换原稿。
\item AI 编剧：入口只在 Word 与 MD。确认可编辑的素材，选择故事构思、人物、大纲、场景、对白、剧集、改编或诊断；建议稿可编辑，确认后插入末尾。若原文已改变，须重新打开工作台再插入。
\end{itemize}
\textbf{MCP 调用}
\begin{itemize}
\item \McpCode{ai_review}：参数 targetId、profileId（默认 science）、language（默认 zh）、literature（默认 true）。返回 source、members、chair、literatureEvidence、partial、inserted=false 和 saved=false。
\item \McpCode{ai_screenwriting}：参数 targetId、task（如 \McpCode{sw-premise-theme}）、medium（电影/短片/电视剧 / 网剧/舞台剧）、instruction。返回 source、content、task、inserted=false 和 saved=false。完整任务与审稿标准枚举见客户端 schema。
\item 目标必须是 Word 或 MD；读取当前整篇未保存正文，超过 120000 字符报错。Word 包含可用对象信息和图片；MD 最多提供 5 张可解析图片。文献或图片未检查时披露限制，不宣称已验证。
\end{itemize}
\textbf{审阅后写回}\par
source 包含目标、路径、脏状态及原文版本。Word 将 source.revision 传给 \McpCode{word_insert_text} 的 expectedRevision；MD 将 source.expectedText 传给 \McpCode{markdown_replace} 并合并原文。版本不符时先重新读取，不覆盖新修改。\par
\Tip{编剧技能来自固定版本的 screenwriting-skills，首次使用需联网。AI 使用桌面配置并可能计费；审稿启用文献检索时会发送检索词。取消会传给模型与文献请求。MCP 不自动写回或保存，界面和接口均保留用户审阅步骤。}
}

\ENPage{47}{AI 审稿与编剧：统一流程}{AI review and screenwriting: shared workflow}{Word, Markdown, and MCP share model settings, screenwriting skills, and the review committee runner. Confirm the source and task, review the result, then write back.}{
\textbf{In the interface}
\begin{itemize}
\item AI review: choose a profile, report language, and literature search. Inspect independent reviews, chair synthesis, and failures. Reports do not automatically replace the draft.
\item Screenwriting belongs to Word and Markdown. Confirm editable source material; choose premise, characters, outline, scene, dialogue, series, adaptation, or diagnosis. Edit the proposal, then append it. If the source document changed, reopen the studio before insertion.
\end{itemize}
\textbf{MCP calls}
\begin{itemize}
\item \McpCode{ai_review}: targetId, profileId (default science), language (default zh), literature (default true). Returns source, members, chair, literatureEvidence, partial, inserted=false, saved=false.
\item \McpCode{ai_screenwriting}: targetId, task (for example \McpCode{sw-premise-theme}), medium, instruction. Returns source, content, task, inserted=false, saved=false. Consult the tool schema for exact task, medium, and review-profile values.
\item Only Word and Markdown targets are accepted. MCP reads the whole current unsaved body and rejects documents above 120000 characters. Word supplies available objects/images; Markdown includes up to five readable images. Missing visual or literature evidence is disclosed.
\end{itemize}
\textbf{Write back after review}\par
source contains the target, path, dirty state, and source version. Pass source.revision as expectedRevision to \McpCode{word_insert_text}; for Markdown, merge the original text and pass source.expectedText to \McpCode{markdown_replace}. Read and merge again if the source changed.\par
\Tip{Screenwriting loads a pinned screenwriting-skills revision and requires network access on first use. Model calls may incur charges; literature search sends search terms. Cancellation propagates to model and literature requests. Neither MCP service automatically inserts or saves its output.}
}
